\chapter{Introduction}
\label{ch:preface}

\section{Objective}
\label{sec:objective}

Everything in the following chapters serves the main research goal of this thesis, which is:

\begin{highlight}
    develop program induction technology for health services based on deep learning
\end{highlight}

\section{Outline and contributions}

Chapter \ref{ch:promise} explains the background and motivation behind the research: what is deep program induction and why is it a promising approach in safety-critical fields like healthcare. Chapter \ref{ch:proposal} proposes a reference architecture to apply deep program induction in healthcare: \emph{Patient Simulator Reinforcement Learning from Code Execution Feedback} (RLCEPS).

RLCEPS requires 2 ingredients: a patient simulator and an algorithm for Reinforcement Learning from Code Execution Feedback.
We review the literature in both areas (chapters \ref{ch:ps-methods} and \ref{ch:simulators-sota}), and argue that the state of the art in each is inadequate to make RLCEPS a reality.
In parts \ref{part:proginduction} and \ref{part:health} we address this with novel contributions to both fields.

Part \ref{part:proginduction} is dedicated to program induction methodology.
Chapter \ref{ch:ps-methods} reviews the state of the art and identifies gaps of particular importance to RLCEPS, particularly the nascent state of \emph{Reinforcement Learning from Code Execution Feedback} (RLCEF).
Chapters \ref{ch:bfpp}, \ref{ch:neurogen}, \ref{ch:tree2tree}, \ref{ch:boptest} and \ref{ch:seidr} propose, evaluate and compare different techniques to address these gaps.
Chapter \ref{ch:bfpp} proposes \emph{BF++}, a novel domain-specific language specifically for RLCEF settings and evaluates an LSTM-based BF++ program synthesis model on standard Gymnasium \cite{towersGymnasiumStandardInterface2024} benchmarks.
Chapter \ref{ch:neurogen} improves upon the results of chapter \ref{ch:bfpp} by combining the strength of the LSTM and genetic programming techniques, a hybrid methodology we refer to as \emph{neurogenetic programming}.
Chapter \ref{ch:tree2tree} explores an evaluates an alternative approach: training a tree-based autoencoder model for \emph{autoencoder genetic programming} on competitive programming data.
Finally, chapters \ref{ch:boptest} and \ref{ch:seidr} explore Reinforcement Learning from Code Execution Feedback with Large Language Models: chapter \ref{ch:boptest} argues that the key to successful application of LLMs in program induction is an iterative process of program improvement and demonstrates this in BOPTEST building optimization benchmark, while chapter \ref{ch:seidr} takes an empirical approach to selecting the optimal iterative program induction algorithm with a comprehensive matrix of experiments comparing different resolutions of the \emph{repair-replace tradeoff} on PSB2 and HumanEval datasets.

Part \ref{part:health} is dedicated to patient simulators for reinforcement learning.
Chapter \ref{ch:simulators-sota} reviews existing patient simulators and identifies major issues, particularly that they are based on dangerously small datasets.
Chapter \ref{ch:heartpole} introduces a simple and deliberately clinically inaccurate simulator to be used as a sanity check.
Chapter \ref{ch:auto-als} introduces Auto-ALS: an \emph{anthropodidactic} emergency care simulator based on a training software package for junior healthcare professionals.
Finally, chapter \ref{ch:mimicseq} makes the first step towards a big data driven patient simulator with a sequential dataset for foundation models in intensive care.

Part \ref{part:conclusion} concludes the thesis: chapter \ref{ch:conclusion} discusses the results and the implications thereof for the future of program synthesis in healthcare.
